


\section{Discussion}

In this paper, we start with the question: What characterizes effective reasoning? The S1 paper and subsequent work \citep{muennighoff2025s1, ringel2025learning, jurayj2025your} suggest scaling test-time compute by inserting \textit{wait} token to the generation, lengthening CoTs and encouraging \Review behaviors. Motivated by this, we ask whether \Length and \Review at test-time correlate with correctness at the CoT level. Contrary to the narrative that “longer and more \Review is better,” we find the opposite: shorter CoT and less \rr are associated with higher accuracy, across both math and scientific reasoning. We hypothesize that \Length and \rr metrics are token-level proxies for deeper structural properties of reasoning. To probe those properties, we introduce a new method for extracting a reasoning graph from the CoT and focus on a simple structural measure: \failedstep. \fsf emerges as the strongest predictor in our study, showing significant correlations for all 10 models on both datasets: lower \fsf reliably correlates with higher accuracy.

Correlations alone do not establish mechanism, so we conduct two causal tests: (1) test-time selection using \fsf and (2) targeted editing that removes failed branches from the CoT. Both interventions significantly improve performance, supporting the view that \fsf is not merely predictive but causal. These results suggest that models have the ability to generate a viable path to the correct answer; however, the presence of failed branches biases subsequent exploration and lowers the overall success rate. 

\yunzhen{Overall, our work provides important insights into effective reasoning by identifying \fsf as a robust predictor of performance. Rather than simply scaling token count—a prevailing focus in test-time scaling—our findings suggest that structural quality, specifically controlling failure propagation, may be a more effective approach. As the field increasingly prioritizes test-time over training compute scaling, these results point toward quality-focused strategies (managing failure propagation through context control \citep{liao2025fractured, team2025kimi}) as a promising complement to quantity-based approaches.}

These insights open several avenues for future research, though important limitations remain. All reported correlations are measured at test time. Understanding how training shapes these test-time behaviors and how to induce low \fsf reasoning during generation remains an important direction for future work. We also note that our analysis operates directly on CoTs under the assumption that a given CoT reflects the model’s reasoning and thus helps characterize effective thinking. Assessing the faithfulness of CoTs \citep{lanham2023measuring, chen2025reasoning} is beyond the scope of this work and is left for future study.
